% Generated draft: requires mathematical and source review.
\documentclass[11pt]{article}
\usepackage[margin=25mm]{geometry}
\usepackage{fontspec}
\setmainfont{Latin Modern Roman}
\usepackage{amsmath,amssymb,mathtools,unicode-math}
\setmathfont{Latin Modern Math}
\usepackage{enumitem,needspace,xurl,hyperref}
\setlength{\parindent}{0pt}
\setlength{\parskip}{5pt}
\setlength{\emergencystretch}{3em}
\newcommand{\N}{\mathbb{N}}
\newcommand{\Z}{\mathbb{Z}}
\newcommand{\Q}{\mathbb{Q}}
\newcommand{\R}{\mathbb{R}}
\newcommand{\C}{\mathbb{C}}
\newcommand{\F}{\mathbb{F}}
\newcommand{\A}{\mathbb{A}}
\newcommand{\PP}{\mathbb{P}}
\newcommand{\E}{\mathbb{E}}
\newcommand{\eps}{\varepsilon}
\newcommand{\dd}{\,\mathrm{d}}
\begin{document}
\section*{UnsolvedMath ID: 20001961}
\textbf{Source problem code:} catalogue entry
\subsection{Definitions and mathematical statement}
\textbf{Source scope.} This record is a mixed contextual extract, not a complete new mathematical question. It contains the tail of Lerman's Comment 5.9 referring back to absent Question 5.8, Boyer's tentative Comment 5.10, and the beginning of Section 6. No missing conjecture is reconstructed.

For a cooriented contact manifold $(M,\alpha)$ with a torus action preserving $\alpha$, define
\[
 \langle\mu_\alpha(x),A\rangle=\alpha_x(A^\#_x),
 \qquad \Psi(x,r)=r\mu_\alpha(x),\quad r>0.
\]
The second map is the homogeneous moment map on the symplectization $(M\times\R_{>0},\dd(r\alpha))$, up to the chosen Hamiltonian sign convention. The source reports contact convexity results [Ler02] using methods of [CDM88]: convexity can fail for two-dimensional tori and connectedness can fail for circles, whereas both hold in torus dimension at least three when the zero level is empty; it says zero avoidance holds in the toric case under its conventions. The extract does not specify the precise moment-image cone or connectedness assertion. In particular a normalized contact image and the homogeneous cone are different objects, and convexity of one is not silently substituted for the other.

\textbf{Supplied auxiliary properness statement.} Let $M$ be any nonempty compact manifold, $V\ne\{0\}$ a finite-dimensional normed vector space, and $\mu:M\to V$ continuous. Put $M^\times=M\setminus\mu^{-1}(0)$ and
\[
 \Psi^\times:M^\times\times\R_{>0}\longrightarrow V\setminus\{0\},
 \qquad (x,r)\longmapsto r\mu(x).
\]
The aid asserts the equivalence
\[
 \Psi^\times\text{ is proper}
 \quad\Longleftrightarrow\quad M^\times\text{ is compact}
 \quad\Longleftrightarrow\quad \mu^{-1}(0)\text{ is open in }M.
\]
Proper means inverse images of compact sets are compact. The empty nonzero locus is allowed and yields an empty-domain proper map. If $M$ is connected and $\mu\not\equiv0$, the criterion reduces to $\mu^{-1}(0)=\varnothing$; the nonzero-map hypothesis prevents the empty-domain exception from being lost.

In the zero-free smooth contact case, choose a Euclidean norm on $V$ and put $\rho(x)=\|\mu_\alpha(x)\|$. There are constants $0<c\le\rho\le C$. For $0<a<b<\infty$, define
\[
 W_{a,b}=\{(x,r):a\le r\rho(x)\le b\}.
\]
The aid describes these as compact exact symplectic cobordisms, with
\[
 \frac aC\le r\le\frac bc,\qquad
 W_{a,b}\cong M\times[a,b],\quad
 (x,s)\longmapsto(x,s/\rho(x)).
\]
On the boundary $r\rho=c_0$, $c_0\in\{a,b\}$, the restriction of the Liouville form is $c_0\alpha/\rho$. The vector field $r\partial_r$ is transverse to both boundary components, pointing inward at the lower boundary and outward at the upper one. These annuli exhaust the symplectization as $a$ decreases to zero and $b$ increases without bound. Nevertheless the full homogeneous map $M\times\R_{>0}\to V$ is nonproper: sequences with $r\to0$ can map into compact sets containing zero. Puncturing the target is therefore essential.

\textbf{Other contextual remarks.} Boyer suggests that Morse theory ``through a range'', as in instanton moduli spaces, might apply to contact moment maps; this is tentative and states no theorem. Section 6 is described as an edited version of T.~Hausel's workshop talk, discussing Higgs-bundle moduli on a Riemann surface, Nakajima quiver varieties, and hypertoric manifolds introduced by Bielawski--Dancer. Its informal common presentation is hyperk\"ahler reduction $T^*A\mathbin{///}G$ of cotangent bundles to affine spaces; it does not provide uniform analytic or dimensional hypotheses for those examples. The Section 6.1 heading ``Generators for the cohomology ring of the quotient'' ends with an unfinished cross-reference, so no generator conjecture is supplied here.
\subsection{Short English statement}
This is context linking contact convexity to a later hyperk\"ahler section, with the actual questions cut off. The supplied auxiliary result characterizes properness of a homogeneous map into the punctured target and gives compact exact symplectic annuli when the contact moment map avoids zero. It does not recover the missing questions.
\subsection{Sources}
\begin{itemize}
\item Local UnsolvedMath registry, permanent ID 20001961, source code catalogue entry; canonical \texttt{aim-geometry-notes.json}, notes[298].
\item AIM \emph{Moment maps and surjectivity in various geometries}, tail of Comment 5.9, Comment 5.10 (C.~Boyer), and opening of Section 6 (edited T.~Hausel talk): \url{https://aimath.org/WWN/momentmaps/momentmaps.pdf}.
\item Supplied reference labels [Ler02] and [CDM88], without full entries. Example attributions: Nakajima quiver varieties and Bielawski--Dancer hypertoric manifolds. Question 5.8 and the forthcoming Section 6.1 conjecture are absent.
\item Registry research attempt 1, dated 2026-08-14, labelled context only and not a problem at run time, rigor ``proved'', novelty confidence low. The explicitly machine-generated, non-peer-reviewed aid must be independently verified before citation.
\end{itemize}
\end{document}
